% GENERATED FILE. Edit canonical lessons, then run npm run generate:books.
% canonical-source-hash: fnv1a64:c60cd12a46179c88
% canonical-lessons: LA-C348-oraculum, LA-C348-vates, LA-C348-prodigium, LA-C348-omen, LA-C348-vaticinor

\chapter{Ōrāculum}
\label{ch:la-oraculum}
\hlchaptermodality{348}

\begin{chapteropening}
\textbf{By the end of this chapter:} \emph{I can talk about signs and prophecy: an oracle, a seer, a portent, an omen and to prophesy.}

Complete the last lesson of chapter 348: I can talk about signs and prophecy: an oracle, a seer, a portent, an omen and to prophesy.
\end{chapteropening}

\section[ōrāculum, ōrāculī]{ōrāculum, ōrāculī — an oracle}

\label{lesson:LA-C348-oraculum}



\begin{quote}
Before the new one: say the Latin for a nymph, then the Latin for a fleece.
\end{quote}

\subsection*{You'll want to know: ōrāculum, ōrāculī}
\textbf{ōrāculum, ōrāculī} — \textquotedblleft{}an oracle\textquotedblright{}.

\begin{quote}
Lēgātī ad ōrāculum missī sunt.
\end{quote}

\begin{quote}
\emph{Envoys were sent to the oracle.}
\end{quote}

Neuter, from \emph{ōrō}, 'to speak, to pray' (taught). Both the place where a god answered and the answer itself.

\begin{cousinweb}
\textbf{ōrāculum} is built on \emph{ōrō}, \textquotedblleft{}to speak, to pray\textquotedblright{}, with the ending \emph{-culum} (Lewis \& Short). English \emph{oracle} comes from it through French.
\end{cousinweb}

\begin{grammarlens}[title={What you've built}]
One more word for telling, explaining and describing at B1.
\end{grammarlens}

\subsection*{Your turn}
\begin{itemize}
\raggedright
  \item \emph{Say it:} \emph{ōrāculum, ōrāculī}
  \item \emph{Say it:} \emph{ōrāculum, ōrāculī}, once more
  \item \emph{Say it:} a sentence of your own with \emph{ōrāculum, ōrāculī}
  \item \emph{Recall:} say the Latin for a nymph, then the Latin for a fleece, then say \emph{ōrāculum, ōrāculī} again
\end{itemize}

\begin{tcolorbox}[breakable,colback=teal!4,colframe=teal!35!black,title={Before you move on}]
What does ōrāculum, ōrāculī mean? (\textquotedblleft{}An oracle\textquotedblright{}.) Say it once more.
\end{tcolorbox}
\section[vātēs, vātis]{vātēs, vātis — a seer, a prophet; a poet}

\label{lesson:LA-C348-vates}



\begin{quote}
Before the new one: say the Latin for a fleece, then the Latin for an oracle.
\end{quote}

\subsection*{You'll want to know: vātēs, vātis}
\textbf{vātēs, vātis} — \textquotedblleft{}a seer, a prophet; a poet\textquotedblright{}.

\begin{quote}
Vātēs fāta urbis cecinit.
\end{quote}

\begin{quote}
\emph{The seer sang of the fates of the city.}
\end{quote}

Masculine or feminine, 3rd declension. Poets liked to call themselves \emph{vātēs}, inspired singers.

\begin{cousinweb}
\textbf{vātēs} is inherited; EtymDB traces it to a Proto-Indo-European word for \textquotedblleft{}seer\textquotedblright{}, and Lewis \& Short compare an Old Irish word for a seer.
\end{cousinweb}

\begin{grammarlens}[title={What you've built}]
One more word for telling, explaining and describing at B1.
\end{grammarlens}

\subsection*{Your turn}
\begin{itemize}
\raggedright
  \item \emph{Say it:} \emph{vātēs, vātis}
  \item \emph{Say it:} \emph{vātēs, vātis}, once more
  \item \emph{Say it:} a sentence of your own with \emph{vātēs, vātis}
  \item \emph{Recall:} say the Latin for a fleece, then the Latin for an oracle, then say \emph{vātēs, vātis} again
\end{itemize}

\begin{tcolorbox}[breakable,colback=teal!4,colframe=teal!35!black,title={Before you move on}]
What does vātēs, vātis mean? (\textquotedblleft{}A seer, a prophet; a poet\textquotedblright{}.) Say it once more.
\end{tcolorbox}
\section[prōdigium, prōdigiī]{prōdigium, prōdigiī — a portent, a wonder}

\label{lesson:LA-C348-prodigium}



\begin{quote}
Before the new one: say the Latin for an oracle, then the Latin for a seer.
\end{quote}

\subsection*{You'll want to know: prōdigium, prōdigiī}
\textbf{prōdigium, prōdigiī} — \textquotedblleft{}a portent, a wonder\textquotedblright{}.

\begin{quote}
Lapidibus pluit; hoc prōdigium omnēs terruit.
\end{quote}

\begin{quote}
\emph{It rained stones; this portent frightened everyone.}
\end{quote}

Neuter. Roman historians list the \emph{prōdigia} of each year: a rain of stones, a cow that spoke. English \emph{prodigy}.

\begin{cousinweb}
\textbf{prōdigium} is explained by Lewis \& Short as standing for \emph{*prōdicium}, from \emph{prōdīcō}, \textquotedblleft{}to foretell\textquotedblright{}: a sign that speaks of what will come. English \emph{prodigy} comes from it.
\end{cousinweb}

\begin{grammarlens}[title={What you've built}]
One more word for telling, explaining and describing at B1.
\end{grammarlens}

\subsection*{Your turn}
\begin{itemize}
\raggedright
  \item \emph{Say it:} \emph{prōdigium, prōdigiī}
  \item \emph{Say it:} \emph{prōdigium, prōdigiī}, once more
  \item \emph{Say it:} a sentence of your own with \emph{prōdigium, prōdigiī}
  \item \emph{Recall:} say the Latin for an oracle, then the Latin for a seer, then say \emph{prōdigium, prōdigiī} again
\end{itemize}

\begin{tcolorbox}[breakable,colback=teal!4,colframe=teal!35!black,title={Before you move on}]
What does prōdigium, prōdigiī mean? (\textquotedblleft{}A portent, a wonder\textquotedblright{}.) Say it once more.
\end{tcolorbox}
\section[ōmen, ōminis]{ōmen, ōminis — an omen, a sign}

\label{lesson:LA-C348-omen}



\begin{quote}
Before the new one: say the Latin for a seer, then the Latin for a portent.
\end{quote}

\subsection*{You'll want to know: ōmen, ōminis}
\textbf{ōmen, ōminis} — \textquotedblleft{}an omen, a sign\textquotedblright{}.

\begin{quote}
Hoc ōmen bonum esse putāvērunt.
\end{quote}

\begin{quote}
\emph{They thought that this omen was good.}
\end{quote}

Neuter, 3rd declension. A chance word or sight that seemed to foretell the future.

\begin{cousinweb}
\textbf{ōmen} had an old form \emph{osmen}, says Varro, quoted by Lewis \& Short; its earlier origin is not settled. English \emph{omen} is the same word.
\end{cousinweb}

\begin{grammarlens}[title={What you've built}]
One more word for telling, explaining and describing at B1.
\end{grammarlens}

\subsection*{Your turn}
\begin{itemize}
\raggedright
  \item \emph{Say it:} \emph{ōmen, ōminis}
  \item \emph{Say it:} \emph{ōmen, ōminis}, once more
  \item \emph{Say it:} a sentence of your own with \emph{ōmen, ōminis}
  \item \emph{Recall:} say the Latin for a seer, then the Latin for a portent, then say \emph{ōmen, ōminis} again
\end{itemize}

\begin{tcolorbox}[breakable,colback=teal!4,colframe=teal!35!black,title={Before you move on}]
What does ōmen, ōminis mean? (\textquotedblleft{}An omen, a sign\textquotedblright{}.) Say it once more.
\end{tcolorbox}
\section[vāticinor, vāticinārī]{vāticinor, vāticinārī — to prophesy, to foretell}

\label{lesson:LA-C348-vaticinor}



\begin{quote}
Before the new one: say the Latin for a portent, then the Latin for an omen.
\end{quote}

\subsection*{You'll want to know: vāticinor, vāticinārī}
\textbf{vāticinor, vāticinārī} — \textquotedblleft{}to prophesy, to foretell\textquotedblright{}.

\begin{quote}
Vātēs bellum longum vāticinātus est.
\end{quote}

\begin{quote}
\emph{The seer foretold a long war.}
\end{quote}

\emph{vāticinor, vāticinārī, vāticinātus sum}, a deponent, from \emph{vātēs} + \emph{canō}.

\begin{cousinweb}
\textbf{vāticinor} is built on \emph{vātēs}, \textquotedblleft{}seer\textquotedblright{} (Lewis \& Short), with a second part from \emph{canō}, \textquotedblleft{}to sing\textquotedblright{}: to sing as a seer.
\end{cousinweb}

\begin{grammarlens}[title={What you've built}]
One more word for telling, explaining and describing at B1.
\end{grammarlens}

\subsection*{Your turn}
\begin{itemize}
\raggedright
  \item \emph{Say it:} \emph{vāticinor, vāticinārī}
  \item \emph{Say it:} \emph{vāticinor, vāticinārī}, once more
  \item \emph{Say it:} a sentence of your own with \emph{vāticinor, vāticinārī}
  \item \emph{Recall:} say the Latin for a portent, then the Latin for an omen, then say \emph{vāticinor, vāticinārī} again
\end{itemize}

\begin{tcolorbox}[breakable,colback=teal!4,colframe=teal!35!black,title={Before you move on}]
What does vāticinor, vāticinārī mean? (\textquotedblleft{}To prophesy, to foretell\textquotedblright{}.) Say it once more.
\end{tcolorbox}
